\documentclass[
  doc,
  floatsintext,
  longtable,
  a4paper,
  nolmodern,
  notxfonts,
  notimes,
  12pt,
  colorlinks=true,linkcolor=blue,citecolor=blue,urlcolor=blue]{apa7}

\usepackage{amsmath}
\usepackage{amssymb}

\geometry{inner=1in, outer=1in}
\fancyhfoffset[LE,RO]{0cm}


\usepackage[bidi=default]{babel}
\babelprovide[main,import]{spanish}


% get rid of language-specific shorthands (see #6817):
\let\LanguageShortHands\languageshorthands
\def\languageshorthands#1{}

\RequirePackage{longtable}
\RequirePackage{threeparttablex}

\makeatletter
\renewcommand{\paragraph}{\@startsection{paragraph}{4}{\parindent}%
	{0\baselineskip \@plus 0.2ex \@minus 0.2ex}%
	{-.5em}%
	{\normalfont\normalsize\bfseries\typesectitle}}

\renewcommand{\subparagraph}[1]{\@startsection{subparagraph}{5}{0.5em}%
	{0\baselineskip \@plus 0.2ex \@minus 0.2ex}%
	{-\z@\relax}%
	{\normalfont\normalsize\bfseries\itshape\hspace{\parindent}{#1}\textit{\addperi}}{\relax}}
\makeatother




\usepackage{longtable, booktabs, multirow, multicol, colortbl, hhline, caption, array, float, xpatch}
\usepackage{subcaption}


\renewcommand\thesubfigure{\Alph{subfigure}}
\setcounter{topnumber}{2}
\setcounter{bottomnumber}{2}
\setcounter{totalnumber}{4}
\renewcommand{\topfraction}{0.85}
\renewcommand{\bottomfraction}{0.85}
\renewcommand{\textfraction}{0.15}
\renewcommand{\floatpagefraction}{0.7}

\usepackage{tcolorbox}
\tcbuselibrary{listings,theorems, breakable, skins}
\usepackage{fontawesome5}

\definecolor{quarto-callout-color}{HTML}{909090}
\definecolor{quarto-callout-note-color}{HTML}{0758E5}
\definecolor{quarto-callout-important-color}{HTML}{CC1914}
\definecolor{quarto-callout-warning-color}{HTML}{EB9113}
\definecolor{quarto-callout-tip-color}{HTML}{00A047}
\definecolor{quarto-callout-caution-color}{HTML}{FC5300}
\definecolor{quarto-callout-color-frame}{HTML}{ACACAC}
\definecolor{quarto-callout-note-color-frame}{HTML}{4582EC}
\definecolor{quarto-callout-important-color-frame}{HTML}{D9534F}
\definecolor{quarto-callout-warning-color-frame}{HTML}{F0AD4E}
\definecolor{quarto-callout-tip-color-frame}{HTML}{02B875}
\definecolor{quarto-callout-caution-color-frame}{HTML}{FD7E14}

%\newlength\Oldarrayrulewidth
%\newlength\Oldtabcolsep


\usepackage{hyperref}



\usepackage{color}
\usepackage{fancyvrb}
\newcommand{\VerbBar}{|}
\newcommand{\VERB}{\Verb[commandchars=\\\{\}]}
\DefineVerbatimEnvironment{Highlighting}{Verbatim}{commandchars=\\\{\}}
% Add ',fontsize=\small' for more characters per line
\usepackage{framed}
\definecolor{shadecolor}{RGB}{241,243,245}
\newenvironment{Shaded}{\begin{snugshade}}{\end{snugshade}}
\newcommand{\AlertTok}[1]{\textcolor[rgb]{0.68,0.00,0.00}{#1}}
\newcommand{\AnnotationTok}[1]{\textcolor[rgb]{0.37,0.37,0.37}{#1}}
\newcommand{\AttributeTok}[1]{\textcolor[rgb]{0.40,0.45,0.13}{#1}}
\newcommand{\BaseNTok}[1]{\textcolor[rgb]{0.68,0.00,0.00}{#1}}
\newcommand{\BuiltInTok}[1]{\textcolor[rgb]{0.00,0.23,0.31}{#1}}
\newcommand{\CharTok}[1]{\textcolor[rgb]{0.13,0.47,0.30}{#1}}
\newcommand{\CommentTok}[1]{\textcolor[rgb]{0.37,0.37,0.37}{#1}}
\newcommand{\CommentVarTok}[1]{\textcolor[rgb]{0.37,0.37,0.37}{\textit{#1}}}
\newcommand{\ConstantTok}[1]{\textcolor[rgb]{0.56,0.35,0.01}{#1}}
\newcommand{\ControlFlowTok}[1]{\textcolor[rgb]{0.00,0.23,0.31}{\textbf{#1}}}
\newcommand{\DataTypeTok}[1]{\textcolor[rgb]{0.68,0.00,0.00}{#1}}
\newcommand{\DecValTok}[1]{\textcolor[rgb]{0.68,0.00,0.00}{#1}}
\newcommand{\DocumentationTok}[1]{\textcolor[rgb]{0.37,0.37,0.37}{\textit{#1}}}
\newcommand{\ErrorTok}[1]{\textcolor[rgb]{0.68,0.00,0.00}{#1}}
\newcommand{\ExtensionTok}[1]{\textcolor[rgb]{0.00,0.23,0.31}{#1}}
\newcommand{\FloatTok}[1]{\textcolor[rgb]{0.68,0.00,0.00}{#1}}
\newcommand{\FunctionTok}[1]{\textcolor[rgb]{0.28,0.35,0.67}{#1}}
\newcommand{\ImportTok}[1]{\textcolor[rgb]{0.00,0.46,0.62}{#1}}
\newcommand{\InformationTok}[1]{\textcolor[rgb]{0.37,0.37,0.37}{#1}}
\newcommand{\KeywordTok}[1]{\textcolor[rgb]{0.00,0.23,0.31}{\textbf{#1}}}
\newcommand{\NormalTok}[1]{\textcolor[rgb]{0.00,0.23,0.31}{#1}}
\newcommand{\OperatorTok}[1]{\textcolor[rgb]{0.37,0.37,0.37}{#1}}
\newcommand{\OtherTok}[1]{\textcolor[rgb]{0.00,0.23,0.31}{#1}}
\newcommand{\PreprocessorTok}[1]{\textcolor[rgb]{0.68,0.00,0.00}{#1}}
\newcommand{\RegionMarkerTok}[1]{\textcolor[rgb]{0.00,0.23,0.31}{#1}}
\newcommand{\SpecialCharTok}[1]{\textcolor[rgb]{0.37,0.37,0.37}{#1}}
\newcommand{\SpecialStringTok}[1]{\textcolor[rgb]{0.13,0.47,0.30}{#1}}
\newcommand{\StringTok}[1]{\textcolor[rgb]{0.13,0.47,0.30}{#1}}
\newcommand{\VariableTok}[1]{\textcolor[rgb]{0.07,0.07,0.07}{#1}}
\newcommand{\VerbatimStringTok}[1]{\textcolor[rgb]{0.13,0.47,0.30}{#1}}
\newcommand{\WarningTok}[1]{\textcolor[rgb]{0.37,0.37,0.37}{\textit{#1}}}

\providecommand{\tightlist}{%
  \setlength{\itemsep}{0pt}\setlength{\parskip}{0pt}}
\usepackage{longtable,booktabs,array}
\newcounter{none} % for unnumbered tables
\usepackage{calc} % for calculating minipage widths
% Correct order of tables after \paragraph or \subparagraph
\usepackage{etoolbox}
\makeatletter
\patchcmd\longtable{\par}{\if@noskipsec\mbox{}\fi\par}{}{}
\makeatother
% Allow footnotes in longtable head/foot
\IfFileExists{footnotehyper.sty}{\usepackage{footnotehyper}}{\usepackage{footnote}}
\makesavenoteenv{longtable}

\usepackage{graphicx}
\makeatletter
\newsavebox\pandoc@box
\newcommand*\pandocbounded[1]{% scales image to fit in text height/width
  \sbox\pandoc@box{#1}%
  \Gscale@div\@tempa{\textheight}{\dimexpr\ht\pandoc@box+\dp\pandoc@box\relax}%
  \Gscale@div\@tempb{\linewidth}{\wd\pandoc@box}%
  \ifdim\@tempb\p@<\@tempa\p@\let\@tempa\@tempb\fi% select the smaller of both
  \ifdim\@tempa\p@<\p@\scalebox{\@tempa}{\usebox\pandoc@box}%
  \else\usebox{\pandoc@box}%
  \fi%
}
% Set default figure placement to htbp
\def\fps@figure{htbp}
\makeatother







\usepackage{newtx}

\defaultfontfeatures{Scale=MatchLowercase}
\defaultfontfeatures[\rmfamily]{Ligatures=TeX,Scale=1}





\title{Sentencias}


\shorttitle{SENTENCIAS}


\usepackage{etoolbox}



\ccoppy{\textcopyright~2026}



\author{Edison Achalma}



\affiliation{
{Escuela Profesional de Economía, Universidad Nacional de San Cristóbal
de Huamanga}}




\leftheader{Achalma}

\date{2026-10-20}


\abstract{Este abstract será actualizado una vez que se complete el
contenido final del artículo. }



\authornote{\par{\addORCIDlink{Edison Achalma}{0000-0001-6996-3364}} 
\par{ }
\par{   El autor no tiene conflictos de interés que revelar.    Los
roles de autor se clasificaron utilizando la taxonomía de roles de
colaborador (CRediT; https://credit.niso.org/) de la siguiente
manera:  Edison Achalma:   conceptualización, metodología, análisis
formal, investigación, recursos, curación de
datos, redacción, visualización, supervisión, administración del
proyecto}
\par{La correspondencia relativa a este artículo debe dirigirse a Edison
Achalma, Escuela Profesional de Economía, Universidad Nacional de San
Cristóbal de Huamanga, Portal Independencia N°
57, Ayacucho, AYA 5001, Perú, Email: \href{mailto:elmer.achalma.09@unsch.edu.pe}{elmer.achalma.09@unsch.edu.pe}}
}

\makeatletter
\let\endoldlt\endlongtable
\def\endlongtable{
\hline
\endoldlt
}
\makeatother

\urlstyle{same}



\hypersetup{
  pdftitle={Sentencias},
  pdfauthor={Autor}
}
\makeatletter
\@ifpackageloaded{caption}{}{\usepackage{caption}}
\AtBeginDocument{%
\ifdefined\contentsname
  \renewcommand*\contentsname{Tabla de contenidos}
\else
  \newcommand\contentsname{Tabla de contenidos}
\fi
\ifdefined\listfigurename
  \renewcommand*\listfigurename{Lista de Figuras}
\else
  \newcommand\listfigurename{Lista de Figuras}
\fi
\ifdefined\listtablename
  \renewcommand*\listtablename{Lista de Tablas}
\else
  \newcommand\listtablename{Lista de Tablas}
\fi
\ifdefined\figurename
  \renewcommand*\figurename{Figura}
\else
  \newcommand\figurename{Figura}
\fi
\ifdefined\tablename
  \renewcommand*\tablename{Tabla}
\else
  \newcommand\tablename{Tabla}
\fi
}
\@ifpackageloaded{float}{}{\usepackage{float}}
\floatstyle{ruled}
\@ifundefined{c@chapter}{\newfloat{codelisting}{h}{lop}}{\newfloat{codelisting}{h}{lop}[chapter]}
\floatname{codelisting}{Listado}
\newcommand*\listoflistings{\listof{codelisting}{Lista de Listados}}
\makeatother
\makeatletter
\makeatother
\makeatletter
\@ifpackageloaded{caption}{}{\usepackage{caption}}
\@ifpackageloaded{subcaption}{}{\usepackage{subcaption}}
\makeatother
\makeatletter
\@ifpackageloaded{fontawesome5}{}{\usepackage{fontawesome5}}
\makeatother

% From https://tex.stackexchange.com/a/645996/211326
%%% apa7 doesn't want to add appendix section titles in the toc
%%% let's make it do it
\makeatletter
\xpatchcmd{\appendix}
  {\par}
  {\addcontentsline{toc}{section}{\@currentlabelname}\par}
  {}{}
\makeatother

%% Disable longtable counter
%% https://tex.stackexchange.com/a/248395/211326

\usepackage{etoolbox}

\makeatletter
\patchcmd{\LT@caption}
  {\bgroup}
  {\bgroup\global\LTpatch@captiontrue}
  {}{}
\patchcmd{\longtable}
  {\par}
  {\par\global\LTpatch@captionfalse}
  {}{}
\apptocmd{\endlongtable}
  {\ifLTpatch@caption\else\addtocounter{table}{-1}\fi}
  {}{}
\newif\ifLTpatch@caption
\makeatother

\begin{document}

\maketitle


\hypertarget{toc}{}
\tableofcontents
\newpage
\section[Introduction]{Sentencias}

\setcounter{secnumdepth}{3}

\setlength\LTleft{0pt}




\subsection{🎯 Motivación}\label{motivaciuxf3n}

Una sentencia es la unidad básica de ejecución en Python: es una
instrucción que le dice al intérprete qué hacer. Entender cómo funcionan
las sentencias es fundamental porque todo programa, desde el más simple
hasta el más complejo, está compuesto por ellas. La forma en que Python
organiza y ejecuta estas sentencias define el flujo de tu programa.

\subsection{💡 Intuición}\label{intuiciuxf3n}

Imagina que estás siguiendo una receta de cocina paso a paso. Cada paso
es una sentencia: ``Pela las papas'', ``Corta en cubos'', ``Calienta el
aceite''. Python ejecuta exactamente en ese orden: primero pela, luego
corta, luego calienta. Así funciona un programa: el intérprete lee tus
sentencias de arriba hacia abajo y las ejecuta en secuencia.

\subsection{🧠 Modelo mental}\label{modelo-mental}

En Python, el flujo de ejecución es \textbf{secuencial por defecto}:

\begin{enumerate}
\def\labelenumi{\arabic{enumi}.}
\tightlist
\item
  \textbf{Una sentencia por línea} normalmente: cada línea contiene una
  instrucción completa.
\item
  \textbf{Ejecución ordenada}: se ejecutan en el mismo orden en que
  están escritas (de arriba a abajo).
\item
  \textbf{Sin punto y coma necesario}: a diferencia de Java o C, Python
  no requiere \texttt{;} para terminar una sentencia; el final de línea
  es suficiente.
\item
  \textbf{Sentencias compuestas}: estructuras como \texttt{if},
  \texttt{for}, \texttt{while} y \texttt{def} ocupan varias líneas y
  definen bloques de código mediante indentación.
\end{enumerate}

\subsection{📐 Desarrollo}\label{desarrollo}

\subsubsection{1. Sentencia básica}\label{sentencia-buxe1sica}

Una sentencia simple es una línea de código que realiza una acción:

\begin{Shaded}
\begin{Highlighting}[]
\BuiltInTok{print}\NormalTok{(}\StringTok{"Hola mundo"}\NormalTok{)  }\CommentTok{\# Sentencia que imprime texto}
\NormalTok{x }\OperatorTok{=} \DecValTok{5}                \CommentTok{\# Sentencia que asigna un valor}
\NormalTok{y }\OperatorTok{=}\NormalTok{ x }\OperatorTok{+} \DecValTok{3}            \CommentTok{\# Sentencia que calcula y asigna}
\end{Highlighting}
\end{Shaded}

Cada una de estas líneas es una sentencia completa. Python las ejecuta
en orden: primero imprime, luego asigna 5 a \texttt{x}, luego calcula y
asigna 8 a \texttt{y}.

\subsubsection{\texorpdfstring{2. Múltiples sentencias en una línea (con
\texttt{;})}{2. Múltiples sentencias en una línea (con ;)}}\label{muxfaltiples-sentencias-en-una-luxednea-con}

\textbf{Poco recomendado}, pero posible:

\begin{Shaded}
\begin{Highlighting}[]
\BuiltInTok{print}\NormalTok{(}\StringTok{"Hola"}\NormalTok{)}\OperatorTok{;} \BuiltInTok{print}\NormalTok{(}\StringTok{"Mundo"}\NormalTok{)}\OperatorTok{;}\NormalTok{ x }\OperatorTok{=} \DecValTok{5}
\end{Highlighting}
\end{Shaded}

Aunque técnicamente válido, esto \textbf{reduce la legibilidad}. El
estilo Python (PEP 8) recomienda:

\begin{Shaded}
\begin{Highlighting}[]
\BuiltInTok{print}\NormalTok{(}\StringTok{"Hola"}\NormalTok{)}
\BuiltInTok{print}\NormalTok{(}\StringTok{"Mundo"}\NormalTok{)}
\NormalTok{x }\OperatorTok{=} \DecValTok{5}
\end{Highlighting}
\end{Shaded}

\subsubsection{3. Sentencias largas y continuación de
línea}\label{sentencias-largas-y-continuaciuxf3n-de-luxednea}

Si una sentencia es demasiado larga, puedes usar
\texttt{\textbackslash{}} para continuar en la siguiente línea:

\begin{Shaded}
\begin{Highlighting}[]
\NormalTok{total }\OperatorTok{=} \DecValTok{100} \OperatorTok{+} \DecValTok{200} \OperatorTok{+} \DecValTok{300} \OperatorTok{+} \OperatorTok{\textbackslash{}}
        \DecValTok{400} \OperatorTok{+} \DecValTok{500} \OperatorTok{+} \DecValTok{600}
\end{Highlighting}
\end{Shaded}

O mejor, usar paréntesis implícitos:

\begin{Shaded}
\begin{Highlighting}[]
\NormalTok{total }\OperatorTok{=}\NormalTok{ (}\DecValTok{100} \OperatorTok{+} \DecValTok{200} \OperatorTok{+} \DecValTok{300} \OperatorTok{+}
         \DecValTok{400} \OperatorTok{+} \DecValTok{500} \OperatorTok{+} \DecValTok{600}\NormalTok{)}
\end{Highlighting}
\end{Shaded}

\subsubsection{4. Sentencias compuestas
(bloques)}\label{sentencias-compuestas-bloques}

Algunas sentencias ocupan varias líneas y definen bloques:

\begin{Shaded}
\begin{Highlighting}[]
\CommentTok{\# Sentencia compuesta: if{-}else}
\ControlFlowTok{if}\NormalTok{ edad }\OperatorTok{\textgreater{}=} \DecValTok{18}\NormalTok{:}
    \BuiltInTok{print}\NormalTok{(}\StringTok{"Eres mayor de edad"}\NormalTok{)}
    \BuiltInTok{print}\NormalTok{(}\StringTok{"Puedes votar"}\NormalTok{)}
\ControlFlowTok{else}\NormalTok{:}
    \BuiltInTok{print}\NormalTok{(}\StringTok{"Eres menor de edad"}\NormalTok{)}

\CommentTok{\# Sentencia compuesta: función}
\KeywordTok{def}\NormalTok{ saludar():}
    \BuiltInTok{print}\NormalTok{(}\StringTok{"Hola"}\NormalTok{)}
    \BuiltInTok{print}\NormalTok{(}\StringTok{"¿Cómo estás?"}\NormalTok{)}

\CommentTok{\# Sentencia compuesta: bucle}
\ControlFlowTok{for}\NormalTok{ i }\KeywordTok{in} \BuiltInTok{range}\NormalTok{(}\DecValTok{3}\NormalTok{):}
    \BuiltInTok{print}\NormalTok{(i)}
    \BuiltInTok{print}\NormalTok{(}\StringTok{"Siguiente"}\NormalTok{)}
\end{Highlighting}
\end{Shaded}

Cada bloque indentado contiene sentencias que se ejecutan juntas.

\subsubsection{5. Sentencias vacías
(pass)}\label{sentencias-vacuxedas-pass}

El \texttt{pass} es una sentencia que no hace nada. Se usa como marcador
de posición:

\begin{Shaded}
\begin{Highlighting}[]
\CommentTok{\# En desarrollo: aún no implementado}
\KeywordTok{def}\NormalTok{ funcion\_futura():}
    \ControlFlowTok{pass}  \CommentTok{\# Hacer nada por ahora}
\end{Highlighting}
\end{Shaded}

\subsection{🔬 Ejemplos}\label{ejemplos}

\textbf{Mínimo ---} Tres sentencias en su propio orden:

\begin{Shaded}
\begin{Highlighting}[]
\BuiltInTok{print}\NormalTok{(}\StringTok{"Primero"}\NormalTok{)}
\BuiltInTok{print}\NormalTok{(}\StringTok{"Segundo"}\NormalTok{)}
\BuiltInTok{print}\NormalTok{(}\StringTok{"Tercero"}\NormalTok{)}
\CommentTok{\# Salida:}
\CommentTok{\# Primero}
\CommentTok{\# Segundo}
\CommentTok{\# Tercero}
\end{Highlighting}
\end{Shaded}

\textbf{Progresivo ---} Sentencias con cálculos:

\begin{Shaded}
\begin{Highlighting}[]
\NormalTok{base }\OperatorTok{=} \DecValTok{10}          \CommentTok{\# Sentencia 1}
\NormalTok{altura }\OperatorTok{=} \DecValTok{5}         \CommentTok{\# Sentencia 2}
\NormalTok{area }\OperatorTok{=}\NormalTok{ base }\OperatorTok{*}\NormalTok{ altura }\OperatorTok{/} \DecValTok{2}  \CommentTok{\# Sentencia 3}
\BuiltInTok{print}\NormalTok{(}\StringTok{"Área:"}\NormalTok{, area)      }\CommentTok{\# Sentencia 4}
\CommentTok{\# Salida: Área: 25.0}
\end{Highlighting}
\end{Shaded}

\textbf{Real ---} Programa con múltiples tipos de sentencias:

\begin{Shaded}
\begin{Highlighting}[]
\CommentTok{\# Sistema de calificaciones simple}
\NormalTok{nombre }\OperatorTok{=} \StringTok{"Ana"}              \CommentTok{\# Asignación}
\NormalTok{nota }\OperatorTok{=} \DecValTok{85}                   \CommentTok{\# Asignación}

\ControlFlowTok{if}\NormalTok{ nota }\OperatorTok{\textgreater{}=} \DecValTok{60}\NormalTok{:              }\CommentTok{\# Sentencia compuesta}
\NormalTok{    estado }\OperatorTok{=} \StringTok{"Aprobado"}     \CommentTok{\# Asignación dentro del bloque}
    \BuiltInTok{print}\NormalTok{(}\StringTok{"¡Felicidades!"}\NormalTok{)  }\CommentTok{\# Sentencia dentro del bloque}
\ControlFlowTok{else}\NormalTok{:}
\NormalTok{    estado }\OperatorTok{=} \StringTok{"Reprobado"}    \CommentTok{\# Asignación en el otro bloque}
    \BuiltInTok{print}\NormalTok{(}\StringTok{"Estudia más"}\NormalTok{)    }\CommentTok{\# Sentencia en el otro bloque}

\BuiltInTok{print}\NormalTok{(}\SpecialStringTok{f"}\SpecialCharTok{\{}\NormalTok{nombre}\SpecialCharTok{\}}\SpecialStringTok{ está }\SpecialCharTok{\{}\NormalTok{estado}\SpecialCharTok{\}}\SpecialStringTok{"}\NormalTok{)  }\CommentTok{\# Sentencia final}
\end{Highlighting}
\end{Shaded}

\subsection{⚠️ Errores comunes}\label{errores-comunes}

{\def\LTcaptype{none} % do not increment counter
\begin{longtable}[]{@{}
  >{\raggedright\arraybackslash}p{(\linewidth - 4\tabcolsep) * \real{0.2681}}
  >{\raggedright\arraybackslash}p{(\linewidth - 4\tabcolsep) * \real{0.3333}}
  >{\raggedright\arraybackslash}p{(\linewidth - 4\tabcolsep) * \real{0.3986}}@{}}
\toprule\noalign{}
\begin{minipage}[b]{\linewidth}\raggedright
Error
\end{minipage} & \begin{minipage}[b]{\linewidth}\raggedright
Síntoma
\end{minipage} & \begin{minipage}[b]{\linewidth}\raggedright
Cómo detectarlo
\end{minipage} \\
\midrule\noalign{}
\endhead
\bottomrule\noalign{}
\endlastfoot
\textbf{Dos sentencias en la misma línea} &
\texttt{SyntaxError:\ invalid\ syntax} & Asegúrate de que cada sentencia
esté en su propia línea \\
\textbf{Olvidar la indentación} en bloques &
\texttt{IndentationError:\ expected\ an\ indented\ block} & Revisa que
los bloques estén correctamente indentados \\
\textbf{Indentación inconsistente} &
\texttt{IndentationError:\ unindent\ does\ not\ match...} & Usa 4
espacios consistentes en todo el bloque \\
\textbf{Punto y coma innecesario} & No da error, pero es anticonvención
& No uses \texttt{;} a menos que sea estrictamente necesario \\
\end{longtable}
}

\subsection{✅ Buenas prácticas}\label{buenas-pruxe1cticas}

\begin{itemize}
\tightlist
\item
  \textbf{Una sentencia por línea}: mejora la legibilidad y facilita la
  depuración.
\item
  \textbf{Evita el punto y coma}: no es necesario en Python y va contra
  las convenciones.
\item
  \textbf{Usa paréntesis para líneas largas}: en lugar de
  \texttt{\textbackslash{}}, usa paréntesis para continuar líneas.
\item
  \textbf{Agrupa sentencias relacionadas}: usa líneas en blanco para
  separar secciones lógicas.
\item
  \textbf{Sigue PEP 8}: la guía de estilo recomienda un máximo de 79
  caracteres por línea.
\end{itemize}

\subsection{📝 Ejercicio}\label{ejercicio}

\textbf{Básico} --- Reorganiza este código para que sea válido en
Python:

\begin{Shaded}
\begin{Highlighting}[]
\BuiltInTok{print}\NormalTok{(}\StringTok{"Inicio del programa"}\NormalTok{) }\BuiltInTok{print}\NormalTok{(}\StringTok{"Segunda instrucción"}\NormalTok{) x}\OperatorTok{=}\DecValTok{5}\NormalTok{ y}\OperatorTok{=}\DecValTok{10} \BuiltInTok{print}\NormalTok{(x }\OperatorTok{+}\NormalTok{ y) }\BuiltInTok{print}\NormalTok{(}\StringTok{"Fin"}\NormalTok{)}
\end{Highlighting}
\end{Shaded}

\textbf{Requisito:} cada sentencia debe estar en su propia línea, y la
legibilidad debe mejorar.

\begin{center}\rule{0.5\linewidth}{0.5pt}\end{center}

\textbf{Solución (no mires hasta haberlo intentado):}

\begin{Shaded}
\begin{Highlighting}[]
\BuiltInTok{print}\NormalTok{(}\StringTok{"Inicio del programa"}\NormalTok{)}
\BuiltInTok{print}\NormalTok{(}\StringTok{"Segunda instrucción"}\NormalTok{)}
\NormalTok{x }\OperatorTok{=} \DecValTok{5}
\NormalTok{y }\OperatorTok{=} \DecValTok{10}
\BuiltInTok{print}\NormalTok{(x }\OperatorTok{+}\NormalTok{ y)}
\BuiltInTok{print}\NormalTok{(}\StringTok{"Fin"}\NormalTok{)}
\end{Highlighting}
\end{Shaded}

\subsection{🗂️ Síntesis}\label{suxedntesis}

\begin{itemize}
\tightlist
\item
  \textbf{Sentencia = instrucción}: cada línea de código que realiza una
  acción.
\item
  \textbf{Orden secuencial}: se ejecutan de arriba a abajo.
\item
  \textbf{Sin punto y coma}: el salto de línea termina cada sentencia.
\item
  \textbf{Sentencias compuestas}: estructuras como \texttt{if} y
  \texttt{for} abarcan múltiples líneas indentadas.
\item
  \textbf{Buena práctica}: una sentencia por línea para máxima
  legibilidad.
\end{itemize}

\subsection{🔗 Qué sigue}\label{quuxe9-sigue}

\begin{itemize}
\tightlist
\item
  \textbf{Relacionado:} variables y tipos de datos (cómo almacenar datos
  que las sentencias procesan)
\item
  \textbf{Siguiente tema:} Variables y tipos de datos --- aprenderás a
  declarar y manejar diferentes tipos de información en Python.
\end{itemize}

\section{Publicaciones Similares}\label{publicaciones-similares}

Si te interesó este artículo, te recomendamos que explores otros blogs y
recursos relacionados que pueden ampliar tus conocimientos. Aquí te dejo
algunas sugerencias:

\begin{enumerate}
\def\labelenumi{\arabic{enumi}.}
\tightlist
\item
  \href{https://numerus-scriptum.netlify.app/python/2020-06-19-instalacion-de-anaconda/index.pdf}{\faIcon{file-pdf}}
  \href{https://numerus-scriptum.netlify.app/python/2020-06-19-instalacion-de-anaconda}{Instalacion
  De Anaconda}
\item
  \href{https://numerus-scriptum.netlify.app/python/2020-06-20-configurar-entorno-virtual-python-anaconda/index.pdf}{\faIcon{file-pdf}}
  \href{https://numerus-scriptum.netlify.app/python/2020-06-20-configurar-entorno-virtual-python-anaconda}{Configurar
  Entorno Virtual Python Anaconda}
\item
  \href{https://numerus-scriptum.netlify.app/python/2021-04-17-01-introducion-a-la-programacion-con-python/index.pdf}{\faIcon{file-pdf}}
  \href{https://numerus-scriptum.netlify.app/python/2021-04-17-01-introducion-a-la-programacion-con-python}{01
  Introducion A La Programacion Con Python}
\item
  \href{https://numerus-scriptum.netlify.app/python/2021-05-31-02-variables-expresiones-y-statements-con-python/index.pdf}{\faIcon{file-pdf}}
  \href{https://numerus-scriptum.netlify.app/python/2021-05-31-02-variables-expresiones-y-statements-con-python}{02
  Variables Expresiones Y Statements Con Python}
\item
  \href{https://numerus-scriptum.netlify.app/python/2021-06-07-03-objetos-de-python/index.pdf}{\faIcon{file-pdf}}
  \href{https://numerus-scriptum.netlify.app/python/2021-06-07-03-objetos-de-python}{03
  Objetos De Python}
\item
  \href{https://numerus-scriptum.netlify.app/python/2021-06-14-04-ejecucion-condicional-con-python/index.pdf}{\faIcon{file-pdf}}
  \href{https://numerus-scriptum.netlify.app/python/2021-06-14-04-ejecucion-condicional-con-python}{04
  Ejecucion Condicional Con Python}
\item
  \href{https://numerus-scriptum.netlify.app/python/2021-06-21-05-iteraciones-con-python/index.pdf}{\faIcon{file-pdf}}
  \href{https://numerus-scriptum.netlify.app/python/2021-06-21-05-iteraciones-con-python}{05
  Iteraciones Con Python}
\item
  \href{https://numerus-scriptum.netlify.app/python/2021-08-16-06-funciones-con-python/index.pdf}{\faIcon{file-pdf}}
  \href{https://numerus-scriptum.netlify.app/python/2021-08-16-06-funciones-con-python}{06
  Funciones Con Python}
\item
  \href{https://numerus-scriptum.netlify.app/python/2021-08-23-07-dataframes-con-python/index.pdf}{\faIcon{file-pdf}}
  \href{https://numerus-scriptum.netlify.app/python/2021-08-23-07-dataframes-con-python}{07
  Dataframes Con Python}
\item
  \href{https://numerus-scriptum.netlify.app/python/2021-11-29-08-prediccion-y-metrica-de-performance-con-python/index.pdf}{\faIcon{file-pdf}}
  \href{https://numerus-scriptum.netlify.app/python/2021-11-29-08-prediccion-y-metrica-de-performance-con-python}{08
  Prediccion Y Metrica De Performance Con Python}
\item
  \href{https://numerus-scriptum.netlify.app/python/2021-12-06-09-metodos-de-machine-learning-para-clasificacion-con-python/index.pdf}{\faIcon{file-pdf}}
  \href{https://numerus-scriptum.netlify.app/python/2021-12-06-09-metodos-de-machine-learning-para-clasificacion-con-python}{09
  Metodos De Machine Learning Para Clasificacion Con Python}
\item
  \href{https://numerus-scriptum.netlify.app/python/2021-12-13-10-metodos-de-machine-learning-para-regresion-con-python/index.pdf}{\faIcon{file-pdf}}
  \href{https://numerus-scriptum.netlify.app/python/2021-12-13-10-metodos-de-machine-learning-para-regresion-con-python}{10
  Metodos De Machine Learning Para Regresion Con Python}
\item
  \href{https://numerus-scriptum.netlify.app/python/2025-05-10-visualizacion-de-datos-con-python/index.pdf}{\faIcon{file-pdf}}
  \href{https://numerus-scriptum.netlify.app/python/2025-05-10-visualizacion-de-datos-con-python}{Visualizacion
  De Datos Con Python}
\end{enumerate}

Esperamos que encuentres estas publicaciones igualmente interesantes y
útiles. ¡Disfruta de la lectura!






\end{document}
